Esta tesis presenta el diseño e implementación de un prototipo de inteligencia de negocio orientado al riesgo crediticio, que integra procesamiento de datos, predicción multi-horizonte, gestión de experimentos y visualización analítica en una institución de la economía social y solidaria del Ecuador.\\

El prototipo presenta un flujo integrado y automatizado del procesamiento de datos y la actualización de los modelos predictivos, consolida la información en un datamart, y proporciona visualizaciones interactivas, respaldado por el registro sistemático de experimentos que facilita el registro y la trazabilidad de los experimentos.\\

La evaluación comparativa entre modelos de gradient boosting, redes neuronales convolucionales y perceptrones multicapa muestra que, bajo la configuración estudiada, los primeros presentan mejor capacidad discriminativa, aunque su rendimiento se degrada en horizontes temporales extendidos.\\

La contribución principal es la integración técnica de un flujo de extremo a extremo—procesamiento, modelado, registro, almacenamiento analítico y visualización—que habilita el análisis crediticio y la observabilidad de experimentos. El prototipo es funcional en el entorno evaluado; su validación operativa y la evaluación de impacto en las decisiones quedan a cargo de la institución objeto del estudio.\\

\noindent\textbf{Palabras clave:}

Riesgo crediticio; predicción multi-horizonte; LightGBM; redes neuronales; inteligencia de negocio; MLOps; datamart; Apache Airflow.